\chapter{Conclusões e Trabalhos Futuros}
\label{chap:conclusoes-e-trabalhos-futuros}

O principal objetivo deste trabalho era desenvolver um agente capaz de jogar o poquêr \textit{Texas Hold'Em}, utilizando de uma rede neural e algoritmos de predição, sendo assim, o objetivo foi alcançado por meio de uma rede neural utilizando aprendizado por reforço e o algoritmo de predição \textit{Monte Carlo}, o agente se mostrou coerente com o proposto, utilizando métodos de desempate robustos e complexos, enquanto eram evitados gargalos, aplicando o motor de pontuação de base 15, uma resolução matemática concreta e que se mostrou extremamente precisa na categorização de cada jogo possível do \textit{Texas Hold'Em}.

Entretanto, pela natureza continua e variável das apostas do \textit{Texas Hold'Em}, o \textit{Q-Learning} não viabiliza tal variação das apostas, já que tornaria os \textit{Q-Values} ainda mais complexos de serem aplicados, sendo assim foi necessário fazer uma discretização das apostas, não dando liberdade para o agente escolher o tamanho da aposta, limitando o mesmo a apenas cinco ações, sendo elas: \textit{Call/Check}, \textit{Raise} leve, \textit{Raise} forte, \textit{fold} e \textit{All-In}. Além de limitações performáticas matriciais, já que o \textit{Node.js} não é tão eficiente quando se trata de calculo de matrizes complexas. O programa também tem um ciclo de apostas fechado, não possibilitando um aumento cada vez maior das apostas, por exemplo, caso o último agente aposte uma quantia, o primeiro agente não é capaz de cobrir essa aposta, ou até mesmo aumenta-la.

Para trabalhos futuros pode-se sugerir a ideia de fazer a substituição do algoritmo de \gls{DQN} por algoritmos capazes de lidar com ações continuas de apostas, tornando o agente ainda mais próximo da realidade. Também pode ser citado fazer uma evolução do \textit{Monte Carlo} permitindo um processamento feito pela \gls{GPU}, acelerando ainda mais o calculo de \textit{Equity}, além de permitir um aumento exponencial da quantidade de simulações. Também pode ser mencionado a possibilidade de haver um jogador humano jogando contra o agente inteligente, permitindo assim testar ainda mais o agente.